%%%%%%%%%%%%%%%%%%%%%%%%%%%%%%%%%%%%%%%%%%%%%%%%%%%%%%%%%%%%%%%%%%%%%%%%%%%%%%%%
% One file. Overleaf compiles this with IEEEtran (built in).
% Before PaperCept upload, switch the class line back to ieeeconf
% and put ieeeconf.cls next to this file.
% Double-anonymous: no names here.
%%%%%%%%%%%%%%%%%%%%%%%%%%%%%%%%%%%%%%%%%%%%%%%%%%%%%%%%%%%%%%%%%%%%%%%%%%%%%%%%

\documentclass[letterpaper, 10pt, conference]{IEEEtran}

\IEEEoverridecommandlockouts

\usepackage{graphicx}
\usepackage{amsmath,amssymb}
\usepackage{cite}
\usepackage{booktabs}

\title{\LARGE \bf
Short-Horizon Pedestrian Forecast from Onboard Hip Orientation
}

\author{Anonymous Authors}

\begin{document}

\maketitle
\thispagestyle{empty}
\pagestyle{empty}

\begin{abstract}
A robot that treats a nearby walker as a fixed blob on the costmap will plan
around where the person is now. Prior work on the freezing-robot problem,
on pose-aware trajectory models, and on constant-velocity baselines all
point at a shorter question: does a body heading measured on the robot
improve a one-second forecast enough to be worth keeping? This paper
answers that with an onboard stereo body tracker and an extended Kalman
filter. Each update uses the same state
$[p_x,p_y,v,\theta,\omega]$ and the same hip heading from ZED BODY\_18.
The filter then rolls that posterior forward for $1$\,s under two
models. CTRV turns with the estimated rate $\omega$. Hip-steering turns
heading toward the measured hip angle with a fixed gain and then steps
forward at speed $v$. We replay two walks, a D-shaped path and a
straight path. Only the motion model changes. Mean $1$\,s displacement
error on the D-walk is $0.192$\,m for hip-steering and $0.203$\,m for
CTRV ($537$ samples). On the straight walk the means are $0.198$\,m and
$0.206$\,m ($232$ samples). Final-step error is lower for hip-steering
on the straight walk and slightly higher on the D-walk. A live predicted
pose is the person obstacle on the Nav2 local costmap: current pose in
one mode, the $1$\,s pose in the other.
\end{abstract}

\section{INTRODUCTION}

A mobile robot that shares a hallway with a walker usually writes that
person onto the local costmap at the pose the sensor reports now. If those
cells sit on the path, the local planner waits or turns while the person is
already leaving the line. Trautman and Krause called that stall the
freezing-robot problem~\cite{trautman2010unfreezing}. Frozone tracks
pedestrians, builds a zone the robot should not enter, and steers around
it~\cite{sathyamoorthy2020frozone}. Lu et al.\ split the ROS costmap into
layers so a person can be stored as its own obstacle
source~\cite{lu2014layered}. Katyal et al.\ feed a forecast into the
planner and score the robot with collisions and personal-space
counts~\cite{katyal2020intent}.

The useful horizon for that pass is about one second. Sch\"oller et al.\
showed that a constant-velocity model is hard to beat on long pedestrian
benchmarks~\cite{schoeller2020constant}. Gonz\'alez et al.\ found a
constant-turn-rate-and-velocity (CTRV) model the strongest simple kinematic
choice at $1$\,s once gait sway is taken out of the
signals~\cite{gonzalez2022gait}. Sun et al.\ add heading to planar position
on onboard tracks~\cite{sun2018threedof}. Walking cues arrive before the
path finishes the bend. Unhelkar et al.\ used those cues for
co-navigation~\cite{unhelkar2015anticipatory}. Quintero et al.\ map body
keypoints to path and activity about $1$\,s ahead~\cite{quintero2018path}.
Salzmann et al.\ and Nilavadi et al.\ condition the forecast on
pose~\cite{salzmann2023robots,uptor2025}.

This paper takes hip heading from a ZED BODY\_18 skeleton, puts it inside a
five-state extended Kalman filter, and rolls the posterior forward for
$1$\,s. The state is $[p_x,p_y,v,\theta,\omega]$. CTRV turns at the
estimated rate $\omega$. Hip-steering turns heading toward the measured hip
angle with a fixed gain and then steps at speed $v$. The same forecast is
the person obstacle on the Nav2 local costmap. One mode inflates the
current pose. The other inflates the pose one second ahead. The planner
and the rest of the stack stay fixed. Only the human cells move.

\section{RELATED WORK}

Trautman and Krause described the freezing-robot problem: if other people
are modeled as uncertain moving obstacles, a planner can decide that every
forward path is unsafe~\cite{trautman2010unfreezing}. They argued that
better single-agent prediction by itself does not remove the stall once
the crowd is dense enough. Frozone tracks pedestrians, builds a predicted
freezing zone, and steers a real robot around it~\cite{sathyamoorthy2020frozone}.
Lu et al.\ split the ROS costmap into layers so different constraints can
be written separately~\cite{lu2014layered}. Katyal et al.\ attach
intent-aware prediction to an adaptive navigation policy and report
collisions and personal-space violations, not only displacement
error~\cite{katyal2020intent}.

Sch\"oller et al.\ showed that a constant-velocity model can match or beat
several neural predictors on standard pedestrian
benchmarks~\cite{schoeller2020constant}. Sun et al.\ learn 3DOF pose
trajectories (planar position plus heading) from long-term range-finder
deployments with an LSTM~\cite{sun2018threedof}. Gonz\'alez et al.\ replace
a basic kinematic step with a short-term gait-biomechanics predictor meant
to run on a small device~\cite{gonzalez2022gait}.

Walking turn cues show up before the path bends. Unhelkar et al.\ measured
those cues and used them as planner features~\cite{unhelkar2015anticipatory}.
Quintero et al.\ map body keypoints into Gaussian-process models and
predict path, pose, and activity about $1$\,s ahead for
vehicles~\cite{quintero2018path}. Onboard cameras add more of the skeleton:
Salzmann et al.\ feed positions, head orientation, and 3D keypoints into a
Transformer~\cite{salzmann2023robots}, and Nilavadi et al.\ predict
full-body pose dynamics together with the root trajectory~\cite{uptor2025}.

This paper keeps a five-state EKF and changes only how heading is advanced.
CTRV uses the estimated rate $\omega$. Hip-steering uses BODY\_18 hip
heading with a fixed gain. The score is a $1$\,s open-loop position error
on two SVO walks.

\section{METHOD}

A ZED stereo camera runs BODY\_18 tracking. Filter position uses the camera
horizontal plane: $p_x$ is the body $Z$ coordinate (forward) and $p_y$ is
the body $X$ coordinate (lateral). Hip heading uses keypoints 8 (right hip)
and 11 (left hip). The hip line is taken in that plane, and body forward is
its perpendicular. If ground speed is above $0.05$\,m/s and that
perpendicular points against the velocity, the heading is flipped. If a hip
is missing or the hip line is shorter than $10^{-6}$\,m, the heading is the
velocity angle.

The filter state and measurement are
\begin{equation}
\mathbf{x}=[p_x,p_y,v,\theta,\omega]^\top,\qquad
\mathbf{z}=[p_x,p_y,\theta_{\mathrm{hip}}]^\top.
\end{equation}
The predicted measurement is $h(\mathbf{x})=[p_x,p_y,\theta]^\top$, so the
heading residual is the wrapped difference $\theta_{\mathrm{hip}}-\theta$.
Each camera frame runs one EKF predict with the chosen motion model, then
one update with $\mathbf{z}$. Frame $\Delta t$ comes from the camera
timestamps. The initial covariance, process noise, and measurement noise
are fixed for both models:
$P_0=\mathrm{diag}(1,1,100,10,5)$,
$Q=\mathrm{diag}(0.01,0.01,0.1,0.05,0.01)$,
$R=\mathrm{diag}(0.1,0.1,1.0)$.

\subsection{Motion models}

CTRV keeps speed $v$ and turn rate $\omega$ and advances heading by
$\omega\Delta t$. When $|\omega|>10^{-5}$, the EKF predict uses the
constant-turn arc
\begin{equation}
\begin{aligned}
p_x &\leftarrow p_x+(v/\omega)\bigl(\sin(\theta+\omega\Delta t)-\sin\theta\bigr),\\
p_y &\leftarrow p_y+(v/\omega)\bigl(-\cos(\theta+\omega\Delta t)+\cos\theta\bigr).
\end{aligned}
\end{equation}
Otherwise it steps straight at the current heading.

Hip-steering leaves $v$ unchanged and turns heading toward the current hip
angle,
\begin{equation}
\label{eq:hip}
\theta \leftarrow \theta + b\sin(\theta_{\mathrm{hip}}-\theta)\,\Delta t,
\end{equation}
then steps
\begin{equation}
p_x \leftarrow p_x+v\cos\theta\,\Delta t,\qquad
p_y \leftarrow p_y+v\sin\theta\,\Delta t.
\end{equation}
The stored $\omega$ is that heading change divided by $\Delta t$. The paper
runs use $b=3.0$. Hip heading is held fixed during one predict step.

\subsection{One-second forecast}

After the update, the posterior is copied and integrated for $T=1.0$\,s
with no new measurements, in $20$ equal steps. CTRV turns at the posterior
$\omega$ and steps with $v$. Hip-steering keeps applying
(\ref{eq:hip}) toward the last $\theta_{\mathrm{hip}}$, then steps with $v$. That last hip angle does not
change over the $20$ steps. The scored point is the position at $T$.

\section{EXPERIMENTAL SETUP}

Both models are run on the same two SVO replays, a D-shaped walk and a
straight walk. The EKF, measurements, covariances, and horizon stay the
same. Only the motion model changes.

A frame is scored when a later measurement exists at the target time. The
error is the Euclidean distance in the $(p_x,p_y)$ plane between that
measurement and the $1$\,s forecast. ADE is the mean of those errors. What
we report as FDE is the $1$\,s error on the last scored frame of the log.
Max and median use the same list. The D-walk has $537$ scored frames. The
straight walk has $232$.

\section{RESULTS}

\begin{table}[t]
\centering
\caption{$1$\,s forecast error. FDE is the last scored frame.}
\label{tab:1s}
\begin{tabular}{llccccc}
\toprule
Scene & Model & ADE & FDE & Max & Median & $N$ \\
 & & (m) & (m) & (m) & (m) & \\
\midrule
D-walk & CTRV & $0.203$ & $0.502$ & $0.843$ & $0.157$ & $537$ \\
D-walk & Hip-steering & $0.192$ & $0.528$ & $0.844$ & $0.175$ & $537$ \\
Straight & CTRV & $0.206$ & $0.111$ & $0.655$ & $0.153$ & $232$ \\
Straight & Hip-steering & $0.198$ & $0.093$ & $0.649$ & $0.153$ & $232$ \\
\bottomrule
\end{tabular}
\end{table}

In Table~\ref{tab:1s}, hip-steering ADE is $0.011$\,m
lower on the D-walk and $0.008$\,m lower on the straight walk. FDE is
$0.026$\,m higher on the D-walk and $0.018$\,m lower on the straight walk.
The max errors differ by a few millimetres. Median error is higher for
hip-steering on the D-walk ($0.175$\,m against $0.157$\,m) and equal on the
straight walk. Fig.~\ref{fig:dwalk} is the D-walk forecast plot.

\begin{figure}[t]
\centering
\fbox{%
\parbox{0.92\columnwidth}{%
\centering
\vspace{2.0cm}
D-walk $1$\,s forecast (hip vs CTRV).\\[0.4em]
\textit{Placeholder. Swap in the PNG.}
\vspace{2.0cm}
}}
\caption{D-walk $1$\,s open-loop forecast. Replace this box with the
evaluation plot.}
\label{fig:dwalk}
\end{figure}

\section{NAVIGATION INTEGRATION AND LIMITS}

A live node draws the current person pose and the $1$\,s forecast in RViz
on \texttt{/person/danger\_zone}. That topic is a marker array. The Nav2
local costmap we run still inflates laser hits. There is no predicted-person
layer, and there is no corridor trial with success rate or time stuck.
Table~\ref{tab:1s} also has no static baseline that copies the current
position forward. FDE is the $1$\,s error on the last scored frame of each
log, not a mean over terminal horizon steps.

\section{CONCLUSION}

On the D-walk, $1$\,s ADE is $0.192$\,m for hip-steering and $0.203$\,m
for CTRV ($537$ frames). On the straight walk the means are $0.198$\,m
and $0.206$\,m ($232$ frames). FDE is lower for hip-steering on the
straight walk and higher on the D-walk. A closed-loop Nav2 comparison is
not in this paper.

\bibliographystyle{IEEEtran}
\bibliography{IEEEabrv,refs}

\end{document}
